\section{Introduction} \label{sec:31-research01-introduction}
\begin{foldable}
  The fruitful performance of \ac{KD} typically relies on the availability of an extensive training set to train the student.
  Traditional \ac{KD} methods often assume that the student is distilled with the same data used for training the teacher~\cite{hinton2015distilling,kim2018paraphrasing,ahn2019variational,tian2020contrastive}.
  However, in real-world scenarios, the distillation often happens at an external party side, where one can only afford a limited-size training set.
  For example, to distill Google's FaceNet (trained using approximately 200 million non-public face images~\cite{schroff2015facenet}) into an in-house model without access to a large-scale dataset like Google's, an external party may only be able to gather a few thousand images or even less.
  This scenario poses the \textit{few-shot} setting and makes it more challenging than standard \ac{KD}.
\end{foldable}

\begin{figure}[ht]
  \centering
  \includegraphics[width=0.99\columnwidth]{./figures/30-research01/31-research01-introduction-standardkd-vs-bbfskd.pdf}
  \caption[Common KD methods and black-box few-shot KD.]{
    Common KD methods and black-box few-shot KD.
    Common KD methods (a) assume no constraints whereas black-box few-shot KD, and (b) only has access to a few images and the teacher's predictive probabilities.
  }
  \label{fig:01-introduction-standard_kd_vs_bbfskd}
\end{figure}

\begin{foldable}
  Another assumption made by traditional \ac{KD} methods is the requirement of a \textit{white-box} teacher \ie, we have complete access to its parameters, activations, or gradients, \etc~\cite{romero2015fitnets,yim2017gift,kimura2018few,chen2019data}.
  However, the teacher model is often \textit{black-box}, \ie, only its final predictions are accessible.
  Understandably, companies often avoid full disclosure of the model parameters as it is their competitive edge and intellectual property.
  For instance, even with a paid subscription to ChatGPT~\cite{openai2022introducing}, only text outputs and their prediction confidence can be retrieved.
  The realistic constraints in the \textit{black-box few-shot} \ac{KD} setting (Fig.~\ref{fig:01-introduction-standard_kd_vs_bbfskd}) render most \ac{KD} methods inapplicable and demand novel approaches.
\end{foldable}

\begin{foldable}
  Recently, some methods have addressed the few-shot \ac{KD} problem, but most of them still require a white-box teacher~\cite{kimura2018few,kong2020learning}.
  To the best of our knowledge, only two methods apply a black-box teacher to few-shot \ac{KD}, namely BBKD~\cite{wang2020neural} and FS-BBT~\cite{nguyen2022black}.
  Both these methods use MixUp~\cite{zhang2018mixup} to generate synthetic images.
  However, as MixUp linearly interpolates between two original images, the synthetic images are either very similar to the original images or do not have realistic semantics, adding little value to the student's training.
  To overcome this problem, FS-BBT employs a conditional variational autoencoder (CVAE)~\cite{sohn2015learning} to generate out-of-distribution synthetic images.
  However, while the pixel-based loss in CVAE is beneficial for reconstructing images, it is not intended for promoting image diversity and have side-effect blurring artifacts~\cite{bond2021deep}.
\end{foldable}

\begin{foldable}
  We address the problem of black-box few-shot \ac{KD} \ie, distilling the student with only a few images and a black-box teacher.
  To boost student performance, we expand the distillation set with diverse synthetic images.
  To achieve diversity in the synthetic images, we propose a novel training scheme for Wasserstein Generative Adversarial Network (WGAN)~\cite{arjovsky2017wasserstein}.
  Firstly, we specifically define \textit{high-confidence} images, which are synthetic images that the teacher can predict confidently w.r.t.\ a threshold, and we use the teacher to compute \textit{adaptive thresholds} in a class-specific and dataset-agnostic manner to avoid biased and abnormal images.
  Then, during the WGAN training loop, we select high-confidence images from the latest iteration on-the-fly and introduce them in the role of additional real images, allowing our WGAN to learn and generate diverse synthetic images close to the teacher's training images.
  Finally, we augment our few-shot images with synthetic images to improve the training of the student, leading to a significantly better performance.
\end{foldable}

\begin{foldable}
  We summarize our contributions as follows:
  \begin{enumerate}
    \item \textbf{DivBFKD}: a method for black-box few-shot \ac{KD} that directly targets diversity of synthetic images as the binding constraint.
    \item A WGAN training scheme that uses the teacher as an external signal, introducing high-confidence synthetic images on-the-fly to anchor generation outside the few-shot information horizon.
    \item Adaptive per-class confidence thresholds that calibrate selection to each class's decision boundary, preventing bias toward easy cases.
    \item Comprehensive evaluation across seven datasets, with ablation confirming that class-wise calibration is load-bearing.
  \end{enumerate}
\end{foldable}

\begin{foldable}
  The chapter is organized as follows.
  \S\ref{sec:32-research01-framework} details the WGAN training scheme and the two mechanisms that anchor synthesis outside the few-shot support: high-confidence image selection and adaptive per-class thresholds.
  \S\ref{sec:33-research01-experiments} evaluates DivBFKD across seven image classification benchmarks and includes ablation studies confirming that class-wise calibration is load-bearing.
  \S\ref{sec:34-research01-conclusion} discusses limitations and connects this contribution to the constraints that motivate Chapter~\ref{ch:40-research02}.
\end{foldable}
